Why does unique factorisation fail for the numbers $a+b\sqrt{-5}$, and come back for their ideals? Why are the carries of school
addition what they are? How many random swaps does it take to shuffle a deck of $52$ cards?

This book is the algebra of the working mathematician. Categories and tensor products; Hilbert's theorems, localisation and
Dedekind domains, where ideals restore unique factorisation; homological algebra and group cohomology, which classify extensions and
explain the carries. Then representations: Wedderburn's theorem, induction, and Burnside's theorem on groups of order $p^aq^b$;
Young diagrams and the hook length formula for the symmetric group; the proof that about $\frac12n\ln n$ random transpositions mix
a deck, with a sharp cutoff; and the Peter--Weyl theorem and Weyl's character formula for the unitary groups.

Every theorem is proved in full, every exercise has a hint and a reasoned solution. Written for the pleasure of seeing why things are true.
